\section*{Asesmen}
\addcontentsline{toc}{section}{Asesmen}

Asesmen Bab V menjadi bagian Tugas Teori (8\%) dan prasyarat
Praktikum 1 (8\%). Fokusnya adalah ketepatan spesifikasi token dan
pemahaman maximal munch.

\begin{table}[H]
\centering
\caption{Instrumen asesmen Bab V}
\label{tab:ch05-asesmen}
\begin{tabularx}{\textwidth}{@{}lXl@{}}
\toprule
\textbf{Instrumen} & \textbf{Indikator} & \textbf{Bobot} \\
\midrule
Tugas Teori & Pattern token, prioritas, dan contoh benar & 8\% \\
P1-A & Spesifikasi token dan 30 kasus uji lengkap & prasyarat \\
Kuis & Token, lexeme, pattern, dan maximal munch & formatif \\
\bottomrule
\end{tabularx}
\end{table}

Kriteria kelulusan: mahasiswa dapat menulis spesifikasi token yang
tidak ambigu, menjelaskan prioritas aturan, dan menghasilkan kasus uji
yang mencakup kasus normal, batas, dan kesalahan.
